\documentclass[11pt]{amsart}
\usepackage[T1]{fontenc}
\usepackage{lmodern}
\usepackage[a4paper,margin=28mm]{geometry}
\usepackage{amsmath,amssymb,amsthm,mathtools}
% Standard typography; no optional microtype processing.
\usepackage{url}
\usepackage[hidelinks]{hyperref}
\newtheorem{theorem}{Theorem}
\newtheorem{lemma}[theorem]{Lemma}
\newcommand{\erfc}{\operatorname{erfc}}
\newcommand{\ee}{\mathrm e}
\newcommand{\dd}{\,\mathrm d}
% Floor delimiters are written explicitly below.
% Continuous equation numbering.
\title[The lane-merging gap near capacity]{A uniform two-term expansion for the\\lane-merging gap near capacity}
\author{Xiaojian Zeng}
\address{PBC School of Finance, Tsinghua University, 100190}
\email{xiaojianzeng@foxmail.com}
\date{September 12, 2026}
\thanks{2020 Mathematics Subject Classification. Primary 05A16; Secondary 60C05, 60F05.}
\keywords{Lane merging, binomial sum, uniform asymptotic expansion, lattice correction}
\hypersetup{pdftitle={A uniform two-term expansion for the lane-merging gap near capacity},pdfauthor={Xiaojian Zeng}}
\begin{document}
\begin{abstract}
We determine the leading term and the constant term of the expected gap
in the limited-capacity lane-merging model when the number of red drivers
is $k=m+\lfloor x\sqrt m\rfloor$. The coefficients are explicit in the
complementary error function and the floor offset. The absolute remainder
is $O_A(m^{-1/2})$, uniformly over every fixed interval $|x|\le A$,
including zero and all floor discontinuities. Binomial symmetry, a local
finite-product estimate, and signed trapezoidal summation give an
elementary proof with explicit, nonoptimized constants. This refines the
near-center approximation question of Insko, Johnson, and Sullivan.
\end{abstract}
\maketitle

\section{The model and the expansion}

In the model of Insko, Johnson, and Sullivan \cite{IJS2026}, cars arrive
one at a time. Red drivers always choose the right lane; green drivers
choose the shorter lane, with ties going to the right. Each lane has
capacity $m$. Once the right lane fills, later cars cannot enter the
left lane. The gap is $m$ minus the number of cars that enter the left
lane. Condition on a uniformly chosen arrival sequence of length $2m-1$
with $k$ red drivers.

Corollary 7.2 of \cite{IJS2026} gives the exact expectation
\begin{equation}\label{eq:definition}
 E_{m,k}=\frac{G_{m,k}}{\binom{2m-1}{k}},\qquad
 G_{m,k}=
 \begin{cases}
 \displaystyle\sum_{i=0}^k\binom{2m-1}{i},&k<m,\\[5pt]
 \displaystyle\sum_{i=0}^{2m-2-k}\binom{2m-1}{i}
 +(k-m+1)\binom{2m}{k+1},&k\ge m,
 \end{cases}
\end{equation}
for integers $m\ge1$ and $0\le k\le2m-1$. A sum with negative upper
endpoint is zero. The upper branch in \cite{IJS2026} includes a final
summand and subtracts $\binom{2m-1}{k}$; binomial symmetry cancels those
two terms and yields \eqref{eq:definition}.
That paper gives off-center upper bounds and a central asymptotic
value, and asks in its Section 8, question 2, for a better approximation
when $k$ is close to $m$. We solve the following precise version:
two explicit terms and a uniform absolute error when $k-m=O(\sqrt m)$.

Use the usual complementary error function
$\erfc(t)=2\pi^{-1/2}\int_t^\infty\ee^{-u^2}\dd u$
\cite[Eq.~7.2.2]{DLMF}. For real $x$ and $-1<\delta\le0$, set
\begin{align}
 H(x)&=\frac{\sqrt\pi}{2}\ee^{x^2}\erfc(|x|)
      =\ee^{x^2}\int_{|x|}^{\infty}\ee^{-t^2}\dd t,
       \qquad x_+=\max(x,0),\label{eq:H}\\
 F_0(x)&=H(x)+2x_+,\qquad
 F_1(x,\delta)=(2\delta+1)xH(x)+\delta+1-2x_+^2.\label{eq:F}
\end{align}

\begin{theorem}\label{thm:main}
For every real $A>0$ there are a finite constant $C_A\ge0$ and an
integer $M_A\ge1$ such that, for every integer $m\ge M_A$ and every
real $|x|\le A$, the index $k=m+\left\lfloor x\sqrt m\right\rfloor$ is admissible and
\begin{equation}\label{eq:main}
 \left|E_{m,k}-\sqrt m\,F_0(x)-F_1(x,\delta_m(x))\right|
 \le \frac{C_A}{\sqrt m},\qquad
 \delta_m(x)=\left\lfloor x\sqrt m\right\rfloor-x\sqrt m.
\end{equation}
Explicit choices of $M_A$ and $C_A$ are given in
\eqref{eq:KM} and \eqref{eq:finalconstants}, with $B=A+1$.
The same constants cover $x=0$, all floor jumps, and $x$ depending on $m$.
\end{theorem}

The earlier unrestricted model of Bardenova, Insko, Johnson, and
Sullivan \cite{BIJS2023} gives a fixed-driver-ratio limit for the
right-lane proportion in Corollary 26 and poses the capacity-gap
expectation separately in Open Problem 39.
Hiveley and Zeilberger \cite{HZ2026} enumerate end states for multiple
lanes in Theorem 1 and tabulate lane-length statistics in Section 6;
their setting has different capacity and conditioning assumptions.
The refinement here is the pair of coefficients in \eqref{eq:F}
together with the uniform error in \eqref{eq:main}.
We found no equivalent expansion in these inspected primary texts
as of September 12, 2026; this comparison is not an exhaustive priority
assessment. The asymptotic tools are classical. We give their needed
inequalities explicitly, including the beta-integral normalization
and trapezoidal estimate; see also \cite[Sections 5.12 and 3.5(i)]{DLMF}.

\section{Local binomial estimates}

Write $h=m^{-1/2}$, $d=k-m$, and $y=dh$. For indices in their supports,
define
\[
 p_{m,j}=2^{-(2m-1)}\binom{2m-1}{m+j},
 \qquad q_{m,j}=4^{-m}\binom{2m}{m+j},
\]
and let $S_{m,d}$ be the distribution function of
$\operatorname{Bin}(2m-1,1/2)$ at $m+d$, extended by zero and one
outside the support. Factorial cancellation gives
\begin{equation}\label{eq:symmetry}
 p_{m,j}=(1-j/m)q_{m,j},\qquad
 q_{m,j}=q_{m,-j},\qquad p_{m,j}=p_{m,-j-1}.
\end{equation}
Thus $S_{m,-1}=1/2$ and
\begin{equation}\label{eq:signedcdf}
 S_{m,d}=
 \begin{cases}
 \displaystyle\frac12+\sum_{j=0}^{d}p_{m,j},&d\ge0,\\[4pt]
 \displaystyle\frac12-\sum_{j=d+1}^{-1}p_{m,j},&d\le-1.
 \end{cases}
\end{equation}
Reflection gives $S_{m,-d-2}=1-S_{m,d}$, with the endpoint $-d-2$
excluding the reflected atom. Consequently \eqref{eq:definition} is
exactly
\begin{equation}\label{eq:branches}
 E_{m,m+d}=
 \begin{cases}
 S_{m,d}/p_{m,d},&d<0,\\[2pt]
 (1-S_{m,d})/p_{m,d}+2m(d+1)/(m+d+1),&d\ge0.
 \end{cases}
\end{equation}
The extra term follows from
$\binom{2m}{m+d+1}/\binom{2m-1}{m+d}=2m/(m+d+1)$.

Fix $B\ge1$ for the remainder of this section, and put
\begin{equation}\label{eq:KM}
 K=\frac14+\frac{5B^2}{6}+\frac{B^4}{3},\qquad
 M=\left\lceil\max\{4,4B^2,K\}\right\rceil,\qquad
 L=\frac32\ee K.
\end{equation}
Also write $f(t)=c\ee^{-t^2}$, where $c=\pi^{-1/2}$.

\begin{lemma}\label{lem:local}
If $m\ge M$, $j$ is an integer with $|j|\le B\sqrt m$, and $t=jh$,
then
\begin{equation}\label{eq:local}
 p_{m,j}=hf(t)(1-th+\rho_{m,j}),\qquad
 |\rho_{m,j}|\le Lh^2,\qquad
 \frac{p_{m,j}}{hf(t)}\ge\frac1{2\ee}.
\end{equation}
\end{lemma}

\begin{proof}
For $I_n=\int_0^{\pi/2}\sin^n u\dd u$, integration by parts gives
$I_n=(n-1)I_{n-2}/n$, with $I_0=\pi/2$ and $I_1=1$. Hence
\[
 I_{2m}=\frac\pi2q_{m,0},\quad
 I_{2m+1}=\frac1{(2m+1)q_{m,0}},\quad
 I_{2m-1}=\frac1{2m q_{m,0}}.
\]
The inequalities $I_{2m+1}\le I_{2m}\le I_{2m-1}$ imply
\begin{equation}\label{eq:wallis}
 \frac1{\sqrt{\pi(m+1/2)}}\le q_{m,0}\le\frac1{\sqrt{\pi m}},
 \qquad
 \left|\log\frac{q_{m,0}}{hf(0)}\right|\le\frac1{4m}.
\end{equation}
The last step uses $\log(1+v)\le v$ for $v\ge0$.

For $r=|j|\le m/2$, use the exact product
\[
 R_{m,r}=\frac{q_{m,r}}{q_{m,0}}
 =\prod_{\ell=1}^{r}\frac{1-(\ell-1)/m}{1+\ell/m}.
\]
For $|u|\le1/2$ the logarithmic series gives
$|\log(1+u)-u+u^2/2|\le 2|u|^3/3$. Applying it to every factor gives
\begin{equation}\label{eq:logproduct}
 \left|\log R_{m,r}+\frac{r^2}{m}-\frac{r^2}{2m^2}\right|
 \le \frac{r^2(r^2+1)}{3m^3}.
\end{equation}
Here the linear sum is $-r^2/m$, the quadratic sum is $r^2/(2m^2)$,
and $\sum_{\ell=1}^r((\ell-1)^3+\ell^3)=r^2(r^2+1)/2$.
Since $r\le B\sqrt m$, equations \eqref{eq:wallis} and
\eqref{eq:logproduct} imply
\[
 q_{m,j}=hf(t)\ee^{\xi_{m,j}},\qquad
 |\xi_{m,j}|\le Kh^2\le1.
\]
Indeed, the noncentral logarithmic error is at most
$(5B^2/6+B^4/3)/m$.
Now $|th|=|j|/m\le1/2$ and $|\ee^\xi-1|\le\ee|\xi|$.
Use \eqref{eq:symmetry} with
$\rho_{m,j}=(1-th)(\ee^{\xi_{m,j}}-1)$ to obtain both assertions
in \eqref{eq:local}.
\end{proof}

The symmetry point $S_{m,-1}=1/2$ lets us estimate the distribution
function using only a finite interval of central masses.
Define
\begin{equation}\label{eq:D}
 D=c\left[(B+1)L+\frac{B(2+4B^2)}{12}
                   +\frac B2+\frac{B(6B+4B^3)}{12}\right],
 \qquad P(y)=\frac12+\int_0^y f(t)\dd t.
\end{equation}

\begin{lemma}\label{lem:cdf}
For $m\ge M$ and integer $|d|\le B\sqrt m$,
\begin{equation}\label{eq:cdf}
 S_{m,d}=P(y)+hf(y)+\varepsilon_{m,d},
 \qquad |\varepsilon_{m,d}|\le Dh^2,\qquad y=dh.
\end{equation}
\end{lemma}

\begin{proof}
For $g\in C^2[-B,B]$, use the signed sum
\[
 \mathcal T_d(g)=
 \begin{cases}
 \displaystyle h\sum_{j=0}^d g(jh),&d\ge0,\\[3pt]
 \displaystyle-h\sum_{j=d+1}^{-1}g(jh),&d\le-1.
 \end{cases}
\]
The trapezoidal estimate, with its orientation retained, is
\begin{equation}\label{eq:trapezoid}
 \left|\mathcal T_d(g)-\int_0^y g(t)\dd t
             -\frac h2(g(0)+g(y))\right|
 \le\frac{|y|h^2}{12}\sup_{|t|\le B}|g''(t)|.
\end{equation}
On one interval, integrating the linear-interpolation error gives
the bound $h^3\sup|g''|/12$, the usual composite trapezoidal remainder
\cite[Eqs.~3.5.2--3.5.3]{DLMF}.
For $d\ge1$, add endpoint half-weights to the trapezoidal sum on
$[0,dh]$. For $d\le-1$, negate the sum on $[dh,0]$ and add the same
half-weights. For $d=0$ the formula is exact. In particular,
\eqref{eq:trapezoid} covers the empty sum at $d=-1$.

Set $g(t)=tf(t)$. Substituting \eqref{eq:local} into
\eqref{eq:signedcdf} yields
\[
 S_{m,d}=\frac12+\mathcal T_d(f)-h\mathcal T_d(g)+\eta,
 \qquad |\eta|\le c(B+1)Lh^2.
\]
There are at most $B/h+1$ terms, each with absolute error at most
$cLh^3$. On $[-B,B]$ we have
\[
 |f''|\le c(2+4B^2),\qquad |g(y)|\le cB,\qquad
 g(0)=0,\qquad |g''|\le c(6B+4B^3).
\]
Apply \eqref{eq:trapezoid} to $f$ and $g$. Since $f'=-2tf$,
\[
 \frac h2(f(0)+f(y))-h\int_0^y tf(t)\dd t=hf(y).
\]
The remaining endpoint term is at most $cBh^2/2$. The quadrature
errors are at most $cB(2+4B^2)h^2/12$ and
$cB(6B+4B^3)h^3/12$. As $h\le1$, their sum with the local errors
is bounded by $Dh^2$, proving \eqref{eq:cdf}.
\end{proof}

\section{Division and the lattice correction}

Continue with $m\ge M$ and $|d|\le B\sqrt m$. Put
\begin{equation}\label{eq:Q}
 \begin{gathered}
 a=c\ee^{-B^2},\qquad J=2\ee(B^2+3L/2),\\
 Q=\frac{J+2\ee D}{a}+B+J,\qquad
 T=4B+2+2(B+1)^3.
 \end{gathered}
\end{equation}
By Lemma~\ref{lem:local},
$z=p_{m,d}/(hf(y))=1-yh+\rho$ is at least $1/(2\ee)$.
The exact reciprocal identity gives
\begin{equation}\label{eq:reciprocal}
 z^{-1}=1+yh+\theta,\qquad
 \theta=\frac{y^2h^2-\rho(1+yh)}{z},\qquad
 |\theta|\le Jh^2.
\end{equation}
The Gaussian integral shows $0\le P(y)\le1$, and $f(y)\ge a$.
Thus \eqref{eq:cdf} and \eqref{eq:reciprocal} imply
\begin{align}
 \left|\frac{S_{m,d}}{p_{m,d}}
 -\left(\frac{P(y)}{hf(y)}+\frac{yP(y)}{f(y)}+1\right)\right|
 &\le Qh,\label{eq:ratio1}\\
 \left|\frac{1-S_{m,d}}{p_{m,d}}
 -\left(\frac{1-P(y)}{hf(y)}+\frac{y(1-P(y))}{f(y)}-1\right)\right|
 &\le Qh.\label{eq:ratio2}
\end{align}
For \eqref{eq:ratio1}, the three errors are bounded by
$(J/a)h$ from $P\theta/(hf)$, $Bh+Jh^2$ from $z^{-1}-1$, and
$(2\ee D/a)h$ from $\varepsilon/(hfz)$. The same bounds prove
\eqref{eq:ratio2} with the sign of the $hf$ correction reversed.
This accounts explicitly for the error from division by a mass of
order $h$.

If $y\ge0$, exact polynomial division, with $v=y+h$, gives
\begin{equation}\label{eq:rational}
 \begin{split}
 \frac{2m(d+1)}{m+d+1}
 &=\frac{2v}{h(1+hv)}\\
 &=\frac{2y}{h}+2-2y^2
       -4yh-2h^2+\frac{2h(y+h)^3}{1+h(y+h)}.
 \end{split}
\end{equation}
The last three terms have absolute value at most $Th$.
For $y<0$, $P(y)/f(y)=H(y)$; for $y\ge0$,
$(1-P(y))/f(y)=H(y)$. Therefore \eqref{eq:branches},
\eqref{eq:ratio1}--\eqref{eq:rational} yield
\begin{equation}\label{eq:grid}
 \left|E_{m,m+d}-h^{-1}F_0(y)-V(y)\right|\le(Q+T)h,
 \qquad V(y)=yH(y)+1-2y_+^2.
\end{equation}
The lower branch includes $d=-1$ and the upper branch includes $d=0$.

\begin{proof}[Proof of Theorem~\ref{thm:main}]
Write $b=\sqrt\pi/2$. The integral
$H(x)=\int_0^\infty\ee^{-u^2-2|x|u}\dd u$ gives $0<H(x)\le b$.
For $x<0$ we have $H'=2xH+1$, and for $x>0$ we have
$H'=2xH-1$. Although $H$ has a corner at zero, $F_0=H+2x_+$
is twice continuously differentiable, with
\begin{equation}\label{eq:smooth}
 F_0'(x)=2xH(x)+1\quad(x\in\mathbb R),\qquad F_0''(0)=2b.
\end{equation}
Indeed, the one-sided first derivatives agree, and for $x\ne0$,
$F_0''=2H+2xH'$ tends to $2b$ at zero. Also $V$ is continuously
differentiable, with $V'(0)=b$ and $V'=H+xH'-4x_+$ away from zero.
Consequently, on $[-B,B]$,
\begin{equation}\label{eq:UW}
 |F_0''|\le U=(2+4B^2)b+2B,\qquad
 |V'|\le W=(1+2B^2)b+5B.
\end{equation}

Given $A>0$, choose
\begin{equation}\label{eq:finalconstants}
 B=A+1,\qquad M_A=M,\qquad C_A=Q+T+U/2+W,
\end{equation}
where \eqref{eq:KM}, \eqref{eq:D}, \eqref{eq:Q}, and \eqref{eq:UW}
define every constant on the right.
For $m\ge M_A$ and $|x|\le A$, let
$d=\left\lfloor x/h\right\rfloor$, $\delta=d-x/h$, and $y=dh=x+\delta h$.
Then $-1<\delta\le0$ and $|y|\le B$. Moreover,
$|d|\le B\sqrt m\le m/2$. Since $m\ge4$, this implies
$0\le m+d\le2m-1$, so the index is admissible.

Taylor's theorem on the interval joining $x$ and $y$ gives
\[
 \left|h^{-1}F_0(y)-h^{-1}F_0(x)-\delta F_0'(x)\right|\le Uh/2,
 \qquad |V(y)-V(x)|\le Wh.
\]
Finally, \eqref{eq:F} and \eqref{eq:smooth} give
$V(x)+\delta F_0'(x)=F_1(x,\delta)$. Combining these inequalities
with \eqref{eq:grid} proves \eqref{eq:main} with the constants in
\eqref{eq:finalconstants}. They depend only on $A$, so the proof
does not restrict the way $x$ varies with $m$ or approaches a floor jump.
\end{proof}

\section{The central cells and finite verification}

The exact central value was already obtained in
\cite[Corollary 7.2]{IJS2026}. From \eqref{eq:branches} one also has
\begin{equation}\label{eq:central}
 E_{m,m-1}=\frac{1}{2q_{m,0}},\qquad
 E_{m,m}-E_{m,m-1}=\frac{m-1}{m+1}.
\end{equation}
At $x=0$, the coefficients are $F_0(0)=\sqrt\pi/2$ and $F_1(0,0)=1$.
In the adjacent negative floor cell, the offset tends to $-1$ as
$x\uparrow0$. This accounts for the distinct constant terms of the two
central cells. More generally, extending \eqref{eq:F} continuously in
$\delta$ to $-1$ gives
\[
 F_1(x,0)-F_1(x,-1)=F_0'(x),
\]
the reset in the lattice correction.

The accompanying program \texttt{supplement/verify.py} checks finite
identities using integers and rational numbers, then performs numerical
consistency checks using mpmath 1.3.0 \cite{mpmath}.
It was written separately from the initial computational checker and
reconstructs the expectations directly from \eqref{eq:definition}.
For $m=1,\ldots,9$ it additionally simulates every binary arrival
sequence using the capacity-blocking rule, and compares the accumulated
gap in each red-count class with $G_{m,k}$. These are $174{,}762$
sequences in $90$ classes. It checks all $4{,}160$ branch identities
for $m=1,\ldots,64$, all $64$ central jumps in that range, and $567$
signed quadrature bounds for monomials of degrees $0,\ldots,6$, meshes
$h=2/7,3/11,5/13$, and $d=-13,\ldots,13$.
Reconstruction with exact predicates and a separate checking program
is also used in the certificate architecture of Zeng, Liu, Ratnavelu,
and Ong \cite{Zeng2026Certificates}, Section 6; the connection here is
the verification procedure.

For the numerical part, let
$A\in\{1/16,3/4,2,4\}$ and take the values of
$m\in\{M_A,M_A+1,3M_A,64,311,1024\}$ satisfying $m\ge M_A$.
Every integer cell intersecting $[-A,A]$ is visited at the phases
\[
 \theta\in
 \left\{0,10^{-30},\frac1{m^2+1},\frac12,
                   1-\frac1{m^2+1},1-10^{-30}\right\},
 \qquad x=\frac{d+\theta}{\sqrt m},
\]
retaining only points in the window. Floor indices and membership are
checked rationally before evaluation. At $80$ decimal digits, the
$1{,}939$ local-mass/CDF comparisons and $7{,}735$ full-expansion
comparisons all satisfy the displayed error bounds.
The special-function computations use floating-point arithmetic,
without interval certification. These finite checks do not exhaust
the real floor cells and do not prove Theorem~\ref{thm:main}.
Its proof is the analytic argument above.

The constants in \eqref{eq:finalconstants} are deliberately conservative;
they establish uniformity rather than a sharp numerical error estimate.
No optimal remainder or theorem for windows growing with $m$ is asserted.

\subsection*{Code and assistance disclosure}
The supplementary program and its recorded output permit reproduction
of the finite checks; its README gives the command and environment.
OpenAI Codex assisted with literature searches, the derivation and
checking of the analytic argument, verification code, and manuscript
preparation. A separate model context reviewed the mathematical dossier;
that review is not outside peer review or statistically independent
evidence. No proof assistant was used. Neither language-model output
nor the numerical checks replace the written proof.

\bibliographystyle{amsplain}
\bibliography{references}
\end{document}
